\chap{Notation and source concordance}\label{app:concordance}
\subsection{Notation}
\begin{longtable}{@{}p{.25\linewidth}p{.7\linewidth}@{}}
\toprule
Symbol & Meaning\\\midrule\endhead
\(\V,F_v,f^*\) & A declared view category, its fibers, and contravariant restrictions.\\
\(\Gamma(F)\) & The strict set of coherent global families.\\
\(\Gamma_U(F),\Ext_U(s)\) & Partial sections on a full patch and their global extension fiber.\\
\(\delta,\partial,d,D\) & Respectively the view differential, tower differential, a model residual, and a Hilbert defect operator. Each target is declared where used.\\
\(p,p_U,p_L\) & A protected observable, restriction to prescribed views, and an innovation observable depending on the splitting.\\
\(\mathfrak o_U,H_U\) & The existence obstruction and the group of invisible coherent changes. Additive and nonabelian conventions are stated separately.\\
\(C_j,J_j,Q_N,S_N\) & Coarsener, chosen lift, restored coarse remainder, and retained finite reconstruction.\\
\(c_f,a_v\) & Coherent nonabelian transition labels and gauge values.\\
\(\mathfrak K_G,J_{\rm comp}\) & The cocycle patch groupoid and its composition-support topology.\\
\(\mathcal E_b(R)\) & Realizations equipped with an actual identification of their boundary with \(b\).\\
\(T_i,K_{ij},R_p\) & Transport to the root, the corresponding root loop operator, and reconstruction from the root.\\
\(\mathcal H,q,P\) & Root residual, an orthogonal kernel projector, and the explicitly defined model repair. A graph repair need not be orthogonal in the global maximum norm.\\
\(\pi_{W,U},H_U\) & The map forgetting part of the prescription and the corresponding invisible coherent ambiguity.\\
\(\phi,\operatorname{EventualHom}\) & A monotone source schedule and transport classes modulo eventual equality.\\
\(\chi_A,M=SA-BC\) & Split exact/residual coordinates and the mismatch between chosen repair lifts.\\
\(\omega,Q_n,C_n,b_n\) & A residual, its prefix-product rate, movement coefficient, and summable step budget. Here $Q_n$ is not the dyadic remainder operator $Q_N$.\\
\(\nu_B,d_W,d_W^{\mathrm{sym}}\) & Residual floor, extended directed transport distance, and its symmetric maximum.\\
\(c(x,y),\kappa(x,y)\) & The two-bit pair phase and its antisymmetric scalar-gauge invariant.\\
\(H,T,r\) & In the budget chapters, $H=\ker p$ with its inherited Hilbert metric, $T=d|_H$, and $r=df$ for the true raw residual. Measured data are declared separately.\\
\(\beta_{T,r}(M),a_\alpha\) & Minimum residual at correction budget $M$, and the quadratic regularized correction.\\
\(C,b,q\) & Closure of the correction image, its projection of the residual, and the irreducible residual floor.\\
\(v,g,e,\mathcal D_\alpha\) & Fit residual, Hilbert normal residual, optimizer error, and regularized dual objective in the solver chapters.\\
\(B,G,q\) & Kernel basis, its Gram matrix, and the coordinate normal covector. This $q$ is not the scalar residual floor used earlier.\\
\(\simeq,\simeq_L\) & A specified equivalence, and a continuous linear equivalence. Neither notation alone asserts an isometry.\\
\bottomrule
\end{longtable}

\subsection{How to read the correspondence}
The identifiers below are stable editorial identifiers, not counts of new discoveries. A row can combine several source declarations into one exposition. The named declaration, its line range, source hash, and correspondence note are recorded in \path{inherited_v1.7.0/data/claim_ledger.json} and in the retained companion Lean guide. Existence of a matching source declaration is not a mechanical proof that every word of an edited statement matches its type. The displayed hypotheses and the printed proof remain part of the review.

\textbf{F} denotes the recovered foundation tree at version 0.20.0. \textbf{M} denotes the focused moving-space tree at version 0.32.0. \textbf{S} denotes the separately recovered Hilbert holonomy source. \textbf{I} denotes a source in the canonical integrated tree, whether inherited or added during the version 1.2.0 formalization. F, M, and S retain provenance, while byte-identical copies of their referenced sources were compiled together in the accepted version 1.2.0 replay of \path{formal/KomorebiSpace/}. All \EditionClaims\ KS catalogue items retain their source correspondences. Their unchanged sources also belong to the accepted version 1.8.0 cumulative scope described below.

The sixteen WR groups are listed separately in Appendix~\ref{app:written}. Nine have complete typed counterparts, two have partial pointwise coverage, and five remain written-only. WR01 and WR02 gained their formal counterparts in version 1.8.0, followed by WR03--WR05 in version 1.10.0. Their declaration map supplements the retained KS catalogue. This does not assert that every illustrative computation or explanatory sentence is independently formalized.

\input{claim_table}

\subsection{Verification is a separate record}
The historical foundation record used a complete build followed by a partial archive replay. The separate moving-space record used two clean builds of its own closure. Neither may be reclassified as a different compilation strategy. The detailed inventories and original receipts are retained in the Lean guide and evidence archives.

The accepted version 1.2.0 record covers \EditionFormalModules\ canonical production modules and a root. Its evidence is bound by \path{inherited_v1.7.0/audit/formal/release.json}. The accepted version 1.6.0 focused record, \path{inherited_v1.7.0/audit/extension/release.json}, covers three then-new modules, three inherited dependencies, and its root. Those historical records retain their original scope.

Version 1.8.0 first rebuilt the canonical core and regularization extension together from an empty project build. All \EditionCumulativeBaseTargets\ source targets passed: \EditionCumulativeBaseModules\ production modules, the two inherited roots, and a combined root. That run audited \EditionCumulativeBaseDeclarations\ project declarations. An independent terminal check verified the complete command inventory, source and control snapshots, and compiled-object hashes.

The continuation then compiled \EditionContinuationModules\ new production modules and the updated combined root over those verified inherited objects. A fresh audit of the entire resulting environment covered \EditionCumulativeDeclarations\ declarations. That version's final source graph has \EditionCumulativeModules\ production modules and \EditionCumulativeRoots\ roots. This was a clean inherited baseline followed by an incremental extension, not a second clean build of all \EditionCumulativeTargets\ final targets. The accepted receipts and source hashes are summarized in \path{continuation/VERIFICATION_SUMMARY.json}. Appendix~\ref{app:written} lists the exact WR coverage and new public declarations.

Version 1.10.0 rebuilds the enlarged project from an empty project build. All \EditionBudgetTargets\ targets comprise \EditionBudgetModules\ production modules and \EditionBudgetRoots\ roots. The four budget modules add \EditionBudgetNewTheorems\ explicit theorem or lemma declarations, bringing that source count to \EditionBudgetTheorems. The full environment receives a fresh declaration-type and axiom audit covering \EditionBudgetDeclarations\ project declarations, including \EditionBudgetNewDeclarations\ from the new modules. All \EditionBudgetControls\ control files meet their specified outcomes: \EditionBudgetPositiveFiles\ positive files, \EditionBudgetCompilerNegativeFiles\ compiler rejection files, and \EditionBudgetPolicyNegativeFiles\ policy rejection files. The accepted run is recorded in
\begin{quote}\small\path{budget_v1.10.0/integration/runs/clean_budget_01/receipt.json}.\end{quote}
Its independent postflight checks bind the source, command inventory, dependency caches, controls, and compiled-object hashes. The precise audited declaration inventory is retained in that record. Only pinned external dependency caches are reused in this clean project replay.

These accepted runs use Lean 4.34.0-rc2 and pinned mathlib commit
\begin{quote}\small\path{aeb72f1e9981b6d35b426b16dca7759be3e7ba06}\end{quote}
The permitted standard axioms are \code{propext}, \code{Classical.choice}, and \code{Quot.sound}. The official compiler, kernel and imported caches remain trusted. Runtime restoration is not a source rebuild of mathlib. Inventory counts include auxiliary or generated declarations as specified by their audit.

Historical outlines retain their dates and scope. Later checked constructions can supersede old targets, as with the composition-site stack theorem. The current declaration, not an early roadmap, determines coverage. Printed proofs, finite tests and typed declarations remain distinct evidence.
